% !TeX root = main.tex
% Engine: XeLaTeX, selected by scripts/build.py.

\documentclass[11pt, aspectratio=169]{beamer}
\input{build/toolkit-path.tex}
\input{\KGUPCToolkitRoot/latex/layouts/beamer.tex}
\usepackage{catchfile}
\usetikzlibrary{arrows.meta,matrix,decorations.pathreplacing}

% 문제 목록의 배치를 유지하고, 아래쪽 빈 공간에 후기 링크를 표시한다.
\newcommand{\afterwordlink}{%
    \begin{tikzpicture}[remember picture,overlay]
        \node[anchor=south,inner sep=0pt,font=\footnotesize]
            at ([yshift=1cm]current page.south)
            {\hyperlink{afterword}{출제진 후기 보기}};
    \end{tikzpicture}%
}

\title[KGUPC 2026 (Advanced) 에디토리얼]
      {KGUPC 2026 (Advanced) \newline 에디토리얼}
\author{Solutio in Kyonggi Univ. KGUPC 2026 출제진}
% 표지의 주최·후원 로고. 추후 스폰서가 너무 많아질 경우, 페이지를 따로 뺄 것을 권장.
\organizerimage{../organizers.png}
\sponsorimage{../sponsors.png}
\date{2026년 10월 30일}

% 문제 제목은 원본 name.tex에서 읽는다. 제목을 중복으로 수정하지 않는다.
\CatchFileDef{\ProblemTitleA}{../problems/A/statement-sections/korean/name.tex}{\endlinechar=-1}
\CatchFileDef{\ProblemTitleB}{../problems/B/statement-sections/korean/name.tex}{\endlinechar=-1}
\CatchFileDef{\ProblemTitleC}{../problems/C/statement-sections/korean/name.tex}{\endlinechar=-1}
\CatchFileDef{\ProblemTitleD}{../problems/D/statement-sections/korean/name.tex}{\endlinechar=-1}
\CatchFileDef{\ProblemTitleE}{../problems/E/statement-sections/korean/name.tex}{\endlinechar=-1}
\CatchFileDef{\ProblemTitleF}{../problems/F/statement-sections/korean/name.tex}{\endlinechar=-1}
\CatchFileDef{\ProblemTitleG}{../problems/G/statement-sections/korean/name.tex}{\endlinechar=-1}
\CatchFileDef{\ProblemTitleH}{../problems/H/statement-sections/korean/name.tex}{\endlinechar=-1}
\CatchFileDef{\ProblemTitleI}{../problems/I/statement-sections/korean/name.tex}{\endlinechar=-1}
\CatchFileDef{\ProblemTitleJ}{../problems/J/statement-sections/korean/name.tex}{\endlinechar=-1}
\CatchFileDef{\ProblemTitleK}{../problems/K/statement-sections/korean/name.tex}{\endlinechar=-1}
\CatchFileDef{\ProblemTitleL}{../problems/L/statement-sections/korean/name.tex}{\endlinechar=-1}

\begin{document}
    \setcounter{framenumber}{-1}
    \frame{\titlepage}

    \begin{frame}{문제 목록}
        \begin{center}
            \begin{tabular}{cl|c|c}
                \hline
                문제 & & 예상 난이도 & 출제자 \\
                \hline
                \problemlink{A}{\textbf{A}} & \problemlink{A}{\ProblemTitleA} & \tiertext{Silver1} & areian13 \\
                \problemlink{B}{\textbf{B}} & \problemlink{B}{\ProblemTitleB} & \tiertext{Silver3} & areian13 \\
                \problemlink{C}{\textbf{C}} & \problemlink{C}{\ProblemTitleC} & \tiertext{Gold3} & yhjang04 \\
                \problemlink{D}{\textbf{D}} & \problemlink{D}{\ProblemTitleD} & \tiertext{Gold1} & ehdnjs14 \\
                \problemlink{E}{\textbf{E}} & \problemlink{E}{\ProblemTitleE} & \tiertext{Gold5} & ehdnjs14 \\
                \problemlink{F}{\textbf{F}} & \problemlink{F}{\ProblemTitleF} & \tiertext{Silver1} & yhjang04 \\
                \problemlink{G}{\textbf{G}} & \problemlink{G}{\ProblemTitleG} & \tiertext{Diamond4} & areian13 \\
                \problemlink{H}{\textbf{H}} & \problemlink{H}{\ProblemTitleH} & \tiertext{Platinum4} & yhjang04 \\
                \problemlink{I}{\textbf{I}} & \problemlink{I}{\ProblemTitleI} & \tiertext{Silver3} & areian13 \\
                \problemlink{J}{\textbf{J}} & \problemlink{J}{\ProblemTitleJ} & \tiertext{Gold2} & yhjang04 \\
                \hline
            \end{tabular}
        \end{center}
    \end{frame}
    \begin{frame}{문제 목록}
        \begin{center}
            \begin{tabular}{cl|c|c}
                \hline
                문제 & & 예상 난이도 & 출제자 \\
                \hline
                \problemlink{K}{\textbf{K}} & \problemlink{K}{\ProblemTitleK} & \tiertext{Silver2} & mingmingmon \\
                \problemlink{L}{\textbf{L}} & \problemlink{L}{\ProblemTitleL} & \tiertext{Platinum5} & areian13 \\
                \hline
            \end{tabular}%
            \afterwordlink
        \end{center}
    \end{frame}

    % 풀이 순서는 문제셋과 동일하다.
    \import{A}{from-npown-to-1.tex}
    \import{B}{cut-dujjonku.tex}
    \import{C}{high-tech-convenience-store.tex}
    \import{D}{connects-grid-pointer.tex}
    \import{E}{ice-water-purifier.tex}
    \import{F}{loop-the-room.tex}
    \import{G}{kaleidoscope-product.tex}
    \import{H}{hongsam-game.tex}
    \import{I}{what-is-my-pair.tex}
    \import{J}{admins-and-reviewers.tex}
    \import{K}{kiryong-mealprep.tex}
    \import{L}{check-check-your-point.tex}
    % !TeX root = main.tex

% 단일 대회에서는 이 파일에 출제진 후기를 작성한다.
% \authorreview{아이디}{후기 본문}을 필요한 만큼 추가한다.
\newcommand{\authorreview}[2]{%
    \par\noindent
    {\color{solutio-purple}\bfseries #1}: #2\par
    \medskip
}

\section{출제진 후기}

% 후기가 길어지면 frame을 추가하고 작성란을 나눈다.
% label=afterword는 첫 후기 프레임에만 둔다.
\begin{frame}[label=afterword]{출제진 후기}
     \authorreview{yhjang04}{
        \textcolor{solutio-gray}{
            이번 KGUPC 다들 많은 관심 가져주셔서 정말 감사드립니다.
        }
    }

    \authorreview{areian13}{
        \textcolor{solutio-gray}{
            후기를 작성해 주세요.
        }
    }

    \authorreview{ehdnjs14}{
        \textcolor{solutio-gray}{
            후기를 작성해 주세요.
        }
    }

    \authorreview{mingmingmon}{
        \textcolor{solutio-gray}{
            제 일상을 녹인 문제들을 출제해서 스토리 짜는데 재밌었습니다. \\
            1회랑 다르게 타대학에서 참가자분들이 많은 관심을 주셔서 대회가 더 풍성해 진 것 같아요. 다음에는 귀여운 특별상도 사비로 기획해오겠습니다 (제 문제 위주로 ㅎㅎ) 
        }
    }


\end{frame}

\end{document}
